% Part: sets-functions-relations
% Chapter: infinite
% Section: dedekinds-proof
%
\documentclass[../../../include/open-logic-section]{subfiles}

\begin{document}

\olfileid[tr]{sfr}{infinite}{dedekindsproof}

\olsection[Dedekind'in ``İspatı'']{Dedekind'in Sonsuz Bir Kümenin
Varlığına İlişkin ``İspatı''}

Bu bölümde doğal sayıları Dedekind cebirleri aracılığıyla küme teorisi
çerçevesinde ele aldık. \olref[arith][ref]{sec}'te, analizin
aritmetikleştirilmesinin (başka şeylerin yanı sıra) felsefi önemi üzerine
düşündük. Şimdi burada başardığımız şeyin önemi üzerine düşünmeliyiz.

\olref[sfr][arith][]{chap} boyunca doğal sayıları verilmiş kabul ettik ve
onları kullanarak tam sayıları, rasyonel sayıları ve gerçel sayıları açıkça
kurduk. Bu bölümdeyse doğal sayıların açık bir kuruluşunu vermedik. Yalnızca,
\emph{herhangi bir Dedekind-sonsuz küme verildiğinde}, $\Nat$'ın davranmasını
istediğimiz gibi davranacak bir küme tanımlayabileceğimizi gösterdik.

Dolayısıyla \emph{hangi nesneler doğal sayılardır} gibi metafizik bir soruyu
yanıtladığımızı ileri süremeyeceğimiz açıktır. Fakat bu iyi bir şeydir.
Nitekim \olref[sfr][arith][ref]{sec}'te, $\Real$'ın Dedekind kesimleri kümesi
olarak tanımını elverişli bir koşul koyma yerine bir \emph{keşif} saymanın
yanlış olacağını vurgulamıştık. Kilit gözlem, Dedekind kesimlerinin gerçel
sayıların temel matematiksel özelliklerini örneklemesidir. Burada da aynı şey
geçerlidir: \emph{her} Dedekind cebiri doğal sayıların temel matematiksel
özelliklerini örnekler. (Gerçekten Dedekind, bütün Dedekind cebirlerinin
\emph{izomorf} olduğunu ispatlayarak bu noktayı pekiştirmiştir
(\citeyear[Theorems 132--3]{Dedekind1888}). Dolayısıyla günümüzdeki birçok
``yapısalcı''nın Dedekind'i bir öncü olarak anması şaşırtıcı değildir.)

Üstelik doğal sayılar teorisinin, kendisi hâlâ oldukça gayriresmî olan ama
doğal sayıları verilmiş kabul etmeyen (öyle görünüyor) naif basit küme
teorisine nasıl gömüleceğini gösterdik. Böylece Dedekind'in kendi iddialı
tasarısını gerçekleştirme yolunda olabiliriz. Dedekind bu tasarıyı şöyle
açıklamıştı:
\begin{quote}
	Bilimde, ispatlanabilir hiçbir şeye ispatı olmaksızın inanılmamalıdır.
	Bu talep akla yatkın görünse de, en basit bilimin temellerini atmaya
	yönelik en yeni yöntemlerde bile; yani mantığın sayılar teorisini ele
	alan bölümünde bile karşılandığını düşünemiyorum. Aritmetikten (cebir,
	analiz) yalnızca mantığın bir parçası olarak söz ederken, sayı kavramını
	uzay ve zaman tasarımlarından ya da sezgilerinden bütünüyle bağımsız
	saydığımı---onu daha ziyade saf düşünce yasalarının dolaysız bir ürünü
	olarak gördüğümü anlatmak istiyorum.
	\citep[preface]{Dedekind1888}
\end{quote}
Dedekind'in cesur düşüncesi şudur. Doğal sayıların yalnızca (naif) küme
teorisi kullanılarak nasıl kurulacağını az önce gösterdik.
\olref[sfr][arith][]{chap}'da doğal sayılar ve bir miktar küme teorisi
verildiğinde gerçel sayıların nasıl kurulacağını gördük. Öyleyse belki de
``aritmetiğin (cebir, analiz)'', Dedekind'in ``mantık'' sözcüğüne verdiği
geniş anlamda, ``yalnızca mantığın bir parçası'' olduğu ortaya çıkar.

Düşünce budur. Fakat bir an duralım. Doğal sayıların yerine kullandığımız
Dedekind cebirini kurmamız, bir Dedekind-sonsuz kümenin varlığı koşuluna
bağlıdır. (\olref[sfr][infinite][dedekind]{thm:DedekindInfiniteAlgebra}'ya
bakmanız yeter.) Bir Dedekind-sonsuz kümenin varlığı ``mantık''la ya da
``saf düşünce yasaları''yla kurulamazsa tasarı burada durur.

Öyleyse bir Dedekind-sonsuz kümenin varlığı ``saf düşünce yasaları''yla
kurulabilir mi? Dedekind'in girişimi şuydu:
\begin{quote}
  Kendi düşünceler alanım, yani düşüncemin nesnesi olabilecek bütün şeylerin
  oluşturduğu $S$ toplamı sonsuzdur. Çünkü $s$, $S$'nin bir elemanını
  gösteriyorsa $s$'nin düşüncemin bir nesnesi olabileceği düşüncesi $s'$ de
  $S$'nin bir elemanıdır. Bunu $s$ elemanının bir görüntüsü $\phi(s)$
  sayarsak, o zaman \dots~$S$ [Dedekind anlamında] sonsuzdur; ispatlanması
  gereken de buydu.
	\citep[\S66]{Dedekind1888}
\end{quote}
Büyük bölümü son derece titiz matematiksel ispatlardan oluşan bir kitabın
ortasında böyle bir şeyle karşılaşmak hayli şaşırtıcıdır. İki noktayı
belirtmek gerekir.

Birincisi: bu ``ispat'', bugün ``matematiksel'' sayacağımız bir nitelik
taşımaktan uzaktır. Psikolojik nesnelerden (düşüncelerden), üstelik yalnızca
\emph{mümkün} olanlardan söz eder.

İkincisi: en azından Dedekind cebirlerini bizim sunduğumuz biçimiyle
alırsak, bu ``ispat''ın açık bir teknik eksiği vardır. Dedekind'in argümanı
başarılıysa yalnızca sonsuz sayıda şey (özellikle sonsuz sayıda düşünce)
bulunduğunu kurar. Fakat Dedekind'in ayrıca $S$'yi, çoğul anlamda
\emph{bazı şeyler} saymak yerine, sonsuz sayıda !!{element} içeren tek bir
\emph{küme} saymak için bize bir gerekçe vermesi gerekir.

Dedekind'in burada bir boşluk görmemiş olması, onun ``toplam'' sözcüğünü
kullanımının \emph{bizim} ``küme'' sözcüğünü kullanımımızla tam olarak
örtüşmediğini düşündürebilir.\footnote{Bunun böyle olduğunu düşünmek için
başka nedenlerimiz de vardır; bkz. \citet[p.~23]{Potter2004}.} Fakat bu pek
şaşırtıcı olmazdı. Son iki bölümde izlediğimiz tasarı---doğal sayıların bir
``kuruluşu'' ve onlardan tam sayıların, gerçel sayıların ve rasyonel sayıların
bir ``kuruluşu''---bütünüyle naifçe yürütüldü. Tam olarak hangi kümelerin ne
zaman var olduğuna ilişkin birçok genel ilke koymadan, ihtiyaç duydukça şu ya
da bu kümeyi hazırdan kullandık. Oysa \emph{bazı} genel ilkelere ihtiyaç
duyduğumuzu biliyoruz; yoksa Russell Paradoksu'na düşeriz.

Naifliğimizi geride bırakmanın zamanı geldi.

\end{document}
